\documentclass[11pt,oneside]{article}
\usepackage[T1]{fontenc}
\usepackage[utf8]{inputenc}
\usepackage{lmodern,microtype}
\usepackage[a4paper,margin=27mm,headheight=15pt]{geometry}
\usepackage{amsmath,amssymb,amsthm,mathtools,mathrsfs}
\usepackage{xcolor,booktabs,longtable,array,enumitem}
\usepackage[most]{tcolorbox}
\usepackage{fancyhdr,needspace}
\usepackage{hyperref}
\usepackage[nameinlink,noabbrev]{cleveref}
\definecolor{ink}{HTML}{253C36}
\definecolor{accent}{HTML}{587A63}
\definecolor{paperbox}{HTML}{F1F4EE}
\hypersetup{colorlinks=true,linkcolor=ink,citecolor=accent,urlcolor=accent,
  pdftitle={Surcomplex Analysis: Hahn-Coherent Holomorphy, Zero Clusters, and Global Rigidity},
  pdfauthor={OpenAI},pdfsubject={Analysis over the complexified surreal numbers}}
\pagestyle{fancy}\fancyhf{}
\fancyhead[L]{\small\sffamily SURCOMPLEX ANALYSIS}
\fancyhead[R]{\small\nouppercase{\leftmark}}
\fancyfoot[C]{\thepage}
\renewcommand{\headrulewidth}{0.3pt}
\renewcommand{\sectionmark}[1]{\markboth{\thesection\quad #1}{}}
\setlength{\parindent}{1.25em}\setlength{\parskip}{3pt}
\setlength{\emergencystretch}{2em}
\allowdisplaybreaks[2]
\numberwithin{equation}{section}
\newtheorem{theorem}{Theorem}[section]
\newtheorem{lemma}[theorem]{Lemma}
\newtheorem{proposition}[theorem]{Proposition}
\newtheorem{corollary}[theorem]{Corollary}
\theoremstyle{definition}
\newtheorem{definition}[theorem]{Definition}
\newtheorem{example}[theorem]{Example}
\newtheorem{construction}[theorem]{Construction}
\theoremstyle{remark}\newtheorem{remark}[theorem]{Remark}
\newcommand{\No}{\mathbf{No}}
\newcommand{\Oz}{\mathbf{Oz}}
\newcommand{\K}{\mathbb K}
\newcommand{\C}{\mathbb C}
\newcommand{\R}{\mathbb R}
\newcommand{\N}{\mathbb N}
\newcommand{\Z}{\mathbb Z}
\newcommand{\OO}{\mathcal O}
\newcommand{\HH}{\mathscr H}
\newcommand{\MM}{\mathscr M}
\newcommand{\m}{\mathfrak m}
\newcommand{\fin}{\mathcal V}
\newcommand{\supp}{\operatorname{supp}}
\newcommand{\st}{\operatorname{st}}
\newcommand{\Res}{\operatorname{Res}}
\newcommand{\ord}{\operatorname{ord}}
\newcommand{\lc}{\operatorname{lc}}
\newcommand{\id}{\operatorname{id}}
\newcommand{\Exp}{\operatorname{Exp}}
\newcommand{\Log}{\operatorname{Log}}
\newcommand{\hoint}{\mathop{\oint}\nolimits^{\!H}}
\newcommand{\hint}{\mathop{\int}\nolimits^{\!H}}
\newcommand{\dd}{\,\mathrm d}
\newcommand{\ev}{\operatorname{ev}}
\newcommand{\Disk}{\mathbb D}
\newcommand{\tube}[1]{#1^{\mu}}
\newtcolorbox{keybox}[1][]{enhanced,breakable,colback=paperbox,colframe=accent,
 boxrule=.5pt,arc=1mm,left=3mm,right=3mm,top=2mm,bottom=2mm,#1}
\newcommand{\term}[1]{\textbf{#1}}
\begin{document}
\begin{titlepage}
\vspace*{16mm}
{\sffamily\small\color{accent} A MATHEMATICAL DEVELOPMENT}\par
\vspace{8mm}
{\Huge\bfseries\color{ink} Surcomplex Analysis}\par
\vspace{5mm}
{\Large Hahn-Coherent Holomorphy,\par Zero Clusters, and Global Rigidity}\par
\vspace{11mm}
{\large Analysis over $\K=\No[i]$}\par
\vspace{5mm}
{\normalsize Prepared for Vladimir Reshetnikov\par 21 September 2026}
\vspace{12mm}
\begin{abstract}
The algebraic closure of the surreal numbers supports a substantial complex
analysis, but not by replacing $\C$ with $\No[i]$ in classical definitions
without further changes. We separate three structures: the full surreal
order topology, infinitesimal formal evaluation, and a coherent analytic
structure obtained from Hahn series with ordinary holomorphic coefficients.
For the coherent structure we prove Taylor and Cauchy formulas, Morera's
criterion, existence and uniqueness of primitives, a residue theorem,
Hahn--Weierstrass preparation, an exact zero-cluster theorem, an argument
principle, and local open and inverse mapping theorems. Infinitesimal
perturbations may involve arbitrary set-sized well-ordered surreal scales,
not just one formal parameter.

We also prove a global rigidity theorem: a function admitting a coherent
Hahn representation on every surreal dilation of a fixed ordinary disk is
a polynomial. We give a normal-form construction of the known canonical
surcomplex exponential, exhibit explicit noncanonical exponentials with
the same local differential equation, and establish a Picard-type
value-distribution theorem for $\Exp(1/z)$. Counterexamples identify exactly
why unrestricted identity, maximum-modulus, and Liouville statements fail.
All infinite sums and contour operations are assigned explicit meanings;
none is silently interpreted as a limit of ordinary partial sums.
\end{abstract}
\vfill
\noindent\small\textbf{Status.} This is a proved framework and research exposition, not a claim
that surcomplex analysis begins here. Surreal normal forms, Hahn summability,
Alling's analytic foundations, and the Ehrlich--Kaplan exponential are prior
work. The structural results below are derived explicitly; no blanket
priority claim or claim of formal verification is made.
\end{titlepage}
\tableofcontents
\newpage

\section{Scope, conventions, and the main results}
\subsection{What a theory has to specify}
A surcomplex number is $x+iy$ with $x,y\in\No$ and $i^2=-1$. The familiar
algebra survives: conjugation, absolute value, polynomial factorization,
and rational differential forms all make sense. Analysis requires choices
that algebra alone does not determine. What counts as an infinite sum?
What domain connects different infinitesimal neighborhoods? Which contours
can be integrated? What conditions exclude a nonconstant locally constant
function?

This article answers these questions with an explicit division of labor.
The \emph{full surreal topology} supplies an honest differential notion.
\emph{Hahn summability} supplies infinite algebraic operations.
\emph{Ordinary holomorphic coefficient functions} supply coherence between
monads with different standard parts. Affine changes of scale place these
constructions near finite, infinite, and infinitesimal surcomplex points.
The resulting theory is neither unrestricted topological calculus nor an
assertion that all classical results transfer automatically.

\subsection{The principal conclusions}
\begin{keybox}
For a connected ordinary complex domain $U$, consider set-supported series
\[
 F=\sum_{\gamma\in S} f_\gamma t^\gamma,
 \qquad f_\gamma\in\OO(U),\qquad t^\gamma:=\omega^{-\gamma},
\]
where $S\subset\No$ is well ordered. They define actual functions on
$U^\mu=\{a+\epsilon:a\in U,\ \epsilon\in\m\}$.

The central results are:
\begin{enumerate}[leftmargin=*,itemsep=2pt]
\item Hahn evaluation is faithful, differentiable, and compatible with
Taylor expansion and coefficientwise Cauchy integration.
\item If the leading coefficient of $F$ has a zero of order $m$ at $a$,
then $F$ has exactly $m$ zeros in $a+\m$, counted with multiplicity.
\item Nonconstant coherent sections are open maps for the full surreal
topology; a nonzero derivative gives a coherent local inverse.
\item Coherence at every surreal scale forces a globally defined function
to be a polynomial.
\end{enumerate}
\end{keybox}
The zero-cluster result is stronger than the assertion that a simple root
persists: it accounts for all roots, including multiple roots splitting at
several widely separated Archimedean scales. Its proof gives a constructive
preparation algorithm. The global result is deliberately a rigidity theorem,
not a proposed definition broad enough to include every transcendental
surcomplex function.

\subsection{Foundations and relation to prior work}
Conway's real-closedness and normal-form theorem, developed further by
Gonshor, are the algebraic starting point \cite{Conway,Gonshor}.
Alling already developed analytic foundations over surreal number fields,
including hyperconvergence and formal implicit-function methods
\cite{Alling,AllingShort}. We do not present those foundational ideas as new.
The extension of restricted analytic functions is also central to the work
of van den Dries and Ehrlich \cite{vdDE}; its summability mechanism is
Neumann's lemma \cite{Neumann}. Our use of holomorphic Hahn coefficients
makes this mechanism explicit and adds a common ordinary domain on which
the coefficients must be coherent.

Peterzil and Starchenko develop a different, powerful framework by requiring
definability in an o-minimal expansion of a real-closed field
\cite{PS,PSsurvey}. Definable connectedness can then replace ordinary
connectedness. We do not identify their definable functions with the Hahn
sections used here. Rubinstein-Salzedo and Swaminathan study other surreal
notions of limits and integration \cite{RSS}; our summation and contour
conventions are specified independently. The existence and kernel of the
canonical exponential on $\No[i]$ are due to Ehrlich and Kaplan
\cite{EK}. Section~\ref{sec:exponential} gives a direct normal-form description
and proves additional differential consequences within the present setup.

\subsection{Set-theoretic and notational conventions}
We work in a conservative class theory such as NBG with global choice.
The symbols $\No$ and $\K$ denote proper classes, but each number and every
support of a Hahn series is a set. Each section is set-coded by its support
and its ordinary coefficient functions. We never sum a genuinely
proper-class family. Every integer indexing a power series, derivative, or
multiplicity is an \emph{ordinary finite} integer.

The notation $t^\gamma=\omega^{-\gamma}$ refers to Conway's monomial map.
It is not the general exponentiation obtained from the surreal exponential.
Thus $t^{\gamma+\delta}=t^\gamma t^\delta$ by the normal-form convention.
The exponent group is the additive ordered group of surreal numbers.
All domains denoted by $U,D,V,\Omega$ without further qualification are
ordinary subsets of $\C$.

\section{The surcomplex field and its summation calculus}
\subsection{Algebra and absolute value}
Set $\K=\No[i]$. For $z=x+iy$, put
\[
 \overline z=x-iy,\qquad |z|=(x^2+y^2)^{1/2}.
\]
Real-closedness of $\No$ gives the positive square root and implies that
$\K$ is algebraically closed. The usual algebraic proof over a real-closed
field applies to each finite polynomial. In particular, all polynomial
degrees and all factorizations used below are finite.

\begin{proposition}
For $z,w\in\K$, one has
\[
 |zw|=|z||w|,\qquad |z+w|\leq |z|+|w|,\qquad
 |\Re z|,|\Im z|\leq |z|.
\]
Every nonzero $z$ has inverse $\overline z/|z|^2$.
\end{proposition}
\begin{proof}
Multiplicativity follows by multiplying $z\bar z$ and $w\bar w$.
The Cauchy--Schwarz inequality in the ordered field $\No$ follows from
$(xv-yu)^2\geq0$. Apply it to the real and imaginary coordinates, then
square both sides of the triangle inequality. The remaining claims follow
from the same identities.
\end{proof}

\subsection{Complex normal forms}
Combining the real and imaginary Conway normal forms gives the unique
representation
\begin{equation}\label{eq:normal}
 z=\sum_{\gamma\in S} c_\gamma t^\gamma,
 \qquad c_\gamma\in\C\setminus\{0\},
 \qquad S\subset\No\text{ a well-ordered set}.
\end{equation}
Conversely, every such set-supported expression is a surcomplex number.
For $z\ne0$ define
\[
 v(z)=\min S,\qquad \lc(z)=c_{v(z)},\qquad v(0)=+\infty.
\]
Then
\[
 v(zw)=v(z)+v(w),\quad
 v(z+w)\geq\min\{v(z),v(w)\},\quad v(|z|)=v(z).
\]
This is an additive valuation with a proper-class value group. The
statements are assertions about the explicitly given normal forms; they
do not require treating a proper class as a set-valued codomain.

\subsection{Finite parts and monads}
Write
\[
 \fin=\{z\in\K:|z|<n\text{ for some }n\in\N\},\qquad
 \m=\{z\in\K:|z|<1/n\text{ for every }n\geq1\}.
\]
Equivalently, $\fin=\{v\geq0\}$ and $\m=\{v>0\}$, with zero included.
The coefficient at exponent zero defines
\[
 \st:\fin\longrightarrow\C,\qquad
 \fin/\m\cong\C.
\]
Thus every finite surcomplex number is uniquely $a+\epsilon$, with
$a\in\C$ and $\epsilon\in\m$. We call $a+\m$ the \emph{monad} of $a$.
More generally, $b+r\m$ is a scaled monad, for $b\in\K$ and $r\ne0$.

\begin{lemma}\label{lem:valuationbound}
If $v(z)\geq\beta$, then $|z|<t^{\beta-1}$.
\end{lemma}
\begin{proof}
The number $t^{-\beta}z$ is finite and therefore has absolute value less
than the positive infinite number $\omega=t^{-1}$.
\end{proof}
The deliberately generous factor $\omega$ gives a bound independent of any
ordinary real size of a leading coefficient.

\subsection{Summable families}
\begin{definition}
A set-indexed family $(z_j)_{j\in J}$ is \emph{Hahn summable} if
$\bigcup_j\supp(z_j)$ is well ordered and each exponent occurs in the
support of only finitely many $z_j$. Its sum is defined by adding the
finitely many coefficients at each exponent.
\end{definition}
The same definition works with coefficients in any commutative ring,
including a ring of holomorphic functions. A complex numerical series
$\sum 2^{-n}$ is \emph{not} Hahn summable when all its summands are placed
at exponent zero. Its ordinary complex sum may instead be computed first
and used as a single coefficient. This distinction is essential.

We use the following classical support theorem as a foundational input.
\begin{lemma}[Neumann support lemma]\label{lem:neumann}
Let $E$ be a well-ordered set of strictly positive elements of an ordered
abelian group. The set $E^*$ of all finite sums of elements of $E$, including
zero, is well ordered. Each group element is the sum of only finitely many
finite ordered words from $E$. Also, sums of two well-ordered subsets are
well ordered, with only finitely many decompositions of a given element
as a sum of one element from each subset.
\end{lemma}
The finite-word assertion includes all possible word lengths. This is what
justifies summing a geometric series, not merely multiplying two Hahn
series. These facts are commonly called Neumann's lemma; see
\cite{Neumann}. One way to understand the combinatorics is that a
subsequence embedding of words over a well-ordered positive alphabet
cannot decrease their sums. A strictly descending sequence of sums, or
infinitely many distinct words of a fixed sum, is therefore excluded by
the well-quasi-ordering of finite words.

\begin{corollary}\label{cor:formal-evaluation}
For $A(X)=\sum_{n\geq0}a_nX^n\in\C[[X]]$ and $\epsilon\in\m$, the
family $(a_n\epsilon^n)_n$ is Hahn summable. Evaluation is compatible with
addition, multiplication, formal differentiation, and composition
$A(B(X))$ when $B(0)=0$.
\end{corollary}
\begin{proof}
Use $E=\supp(\epsilon)\subset\No^{>0}$. Each term is supported on
sums of $n$ elements of $E$. The support lemma permits every regrouping
in the stated formal identities. At $\epsilon=0$ the assertion is
immediate. The same reasoning with the union of finitely many positive
supports proves multivariable versions.
\end{proof}
No positive complex radius of convergence is required. For instance,
$\sum n!\epsilon^n$ is meaningful for every surreal infinitesimal
$\epsilon$, although the associated complex series has radius zero.

\section{Topology: the obstruction to naive complex analysis}
\subsection{The full surreal topology}
Use balls
\[
 B(a,r)=\{z\in\K:|z-a|<r\},\qquad r\in\No^{>0}.
\]
Both $\fin$ and $\m$, and all their affine translates and nonzero
scalar multiples, are open and closed. For $\m$, a ball of infinitesimal
radius about any of its points stays inside $\m$. A point outside $\m$
has $|z|\geq1/n$ for some finite $n$, so a ball of radius $1/(2n)$ stays
outside. The proof for $\fin$ is analogous, using a ball of radius one
around an infinite point.

\begin{theorem}[No nontrivial set-sized convergence]\label{thm:no-convergence}
Every set of surcomplex numbers is closed and discrete in the full surreal
topology. A net indexed by a set converges in this topology only if it is
eventually equal to its limit. In particular, every convergent ordinary
sequence is eventually constant.
\end{theorem}
\begin{proof}
Fix a set $A\subset\K$ and $p\in\K$. The nonzero distances
$\{|a-p|:a\in A,\ a\ne p\}$ form a set of positive surreal numbers.
The surreal cut property supplies a positive $r$ smaller than every one of
these distances. Then $B(p,r)\cap A$ is empty or equals $\{p\}$.
This proves closedness and discreteness. Apply the same construction to
the set of all nonzero distances between a net's values and its proposed
limit. Eventual membership in the resulting ball forces eventual equality.
\end{proof}
The cut is an application of the simplicity theorem, not an infimum or
an order limit. This does not make the full class $\K$ discrete: every ball
contains other points. Nor does it identify the induced full-surreal topology
on a set-sized Hahn subfield with that subfield's own valuation topology. In particular,
\[
 \sum_{n\geq0}t^n=(1-t)^{-1}
\]
is a Hahn identity, not convergence of finite partial sums in $\K$.
For $t=\omega^{-1}$, the radius $\omega^{-\omega}$ already prevents that
sequence of partial sums from converging to its Hahn sum.

\begin{corollary}\label{cor:no-path}
Every continuous map from the ordinary real interval $[0,1]$ into $\K$
with its full surreal topology is constant.
\end{corollary}
\begin{proof}
Its image is a set and hence a discrete subspace by
\cref{thm:no-convergence}. The continuous image of an ordinary connected
space is connected, so this image is a singleton.
\end{proof}
Consequently, an ordinary parametrized complex circle is not a nonconstant
continuous path in this topology. Coefficientwise contours introduced
later must not be mistaken for such paths.

\subsection{Differentiation still makes sense}
\begin{definition}
A function $F$ on a full-open class domain is \emph{differentially
holomorphic} at $a$ if some $L\in\K$ satisfies
\[
 \forall\eta>0\ \exists\delta>0\ \forall h\in\K:
 \quad 0<|h|<\delta\ \Longrightarrow\
 \left|\frac{F(a+h)-F(a)}h-L\right|<\eta.
\]
The radii $\eta,\delta$ range over all positive surreals. Write $F'(a)=L$.
\end{definition}
Uniqueness, the sum and product rules, and the chain rule follow from the
usual finite inequalities over an ordered field. No limiting sequence is
needed in their proofs. The definition is not vacuous: its quantification
over all sufficiently small increments is a proper-class quantification.

\begin{example}[Bounded entire functions with zero derivative]\label{ex:locallyconstant}
The function
\[
 L(z)=\begin{cases}0,&z\in\fin,\\1,&z\notin\fin\end{cases}
\]
is locally constant and therefore differentially holomorphic on all of
$\K$, with every positive-order derivative zero. It is bounded and
nonconstant. It vanishes on a full-open class, has local maxima of its
absolute value, and is not an open map. Adding $L$ to a primitive preserves
its derivative.
\end{example}
Thus unrestricted analogues of the identity theorem, the strong maximum
principle, Liouville's theorem, and uniqueness of primitives are false.
This is a failure of hypotheses, not a difficulty repaired by a longer
classical proof.

\subsection{The Cauchy--Riemann equations}
Suppose $F=u+iv$ is \emph{real Fr\'echet differentiable} at a point: it has
a real-linear first-order approximation with remainder $o(|h|)$ in the
full surreal sense. Here real-linear means $\No$-linear on the two
surreal-real coordinates, not merely $\R$-linear on $\No$ as an
infinite-dimensional vector space. Then $F$ is differentially holomorphic there exactly
when
\[
 u_x=v_y,\qquad u_y=-v_x.
\]
Indeed, these equations say precisely that the real-linear differential
is multiplication by one element of $\K$. Existence of partial derivatives
alone is not being substituted for Fr\'echet differentiability. In the
coherent category below, Taylor remainders prove that differentiability
directly.

\section{Infinitesimal evaluation and canonical complex lifts}
\subsection{Formal functions on a monad}
For $A\in\C[[X]]$, define $A(\epsilon)$ on $\m$ by
\cref{cor:formal-evaluation}. Formal substitution in two infinitesimals
gives
\begin{equation}\label{eq:formal-taylor}
 A(\epsilon+h)=\sum_{n\geq0}\frac{A^{(n)}(\epsilon)}{n!}h^n,
 \qquad \epsilon,h\in\m.
\end{equation}
After removing the terms through degree $N$, the remainder is
$h^{N+1}Q$, where $v(Q)\geq0$. Thus $|Q|<\omega$ and the formal
derivative is also the full surreal derivative. All higher derivatives
exist. Formal Taylor analyticity on a monad is therefore strictly broader
than ordinary complex analyticity.

The series $\sum n!X^n$ provides a useful test: it is legitimate on $\m$,
but it cannot be the restriction of a coherent section on an ordinary
neighborhood of zero having those same derivatives. Its exponent-zero
coefficient would have Taylor coefficients $n!$, contradicting positive
complex radius of convergence.

\subsection{The canonical lift}
For an ordinary open set $U\subset\C$, define its \emph{tube}
\[
 U^\mu=\{a+\epsilon:a\in U,\ \epsilon\in\m\}.
\]
For $f\in\OO(U)$, set
\begin{equation}\label{eq:lift}
 \widetilde f(a+\epsilon)
   =\sum_{n\geq0}\frac{f^{(n)}(a)}{n!}\epsilon^n.
\end{equation}
The decomposition into standard part and infinitesimal part is unique,
so this is a definition on the whole tube.

\begin{theorem}[Functorial lifting]\label{thm:lift}
The map $f\mapsto\widetilde f$ is an injective algebra homomorphism,
commutes with differentiation and restriction, and satisfies
\[
 \widetilde{g\circ f}=\widetilde g\circ\widetilde f
 \quad\text{whenever }f(U)\subset V,\ g\in\OO(V).
\]
Moreover $\st(\widetilde f(z))=f(\st z)$.
\end{theorem}
\begin{proof}
Evaluation at ordinary points proves injectivity. Taylor coefficients
satisfy the formal product and composition identities. Apply the support
lemma to evaluate those identities at an infinitesimal. The constant
Taylor coefficient is the standard part. Differentiation follows from
\eqref{eq:formal-taylor} and its remainder estimate.
\end{proof}

\begin{corollary}[Lifted local inverse]
If $f'(a)\ne0$, the classical local inverse $g$ lifts to an inverse of
$\widetilde f:a+\m\longrightarrow f(a)+\m$.
\end{corollary}
\begin{proof}
The inverse identities lift by composition. Every point of either monad
belongs to the tube of the corresponding ordinary inverse neighborhood.
\end{proof}
Canonical lifts preserve ordinary analytic continuation because one first
continues the coefficient function and then applies \eqref{eq:lift}. This
is a coherence requirement, not a consequence of monads being connected.

\section{Hahn-coherent holomorphic sections}\label{sec:coherent}
\subsection{Definition and evaluation}
For a connected ordinary domain $U$, let
\begin{equation}\label{eq:Hring}
 \HH(U)=\OO(U)((t^{\No}))_{\mathrm{set}}
 =\left\{\sum_{\gamma\in S}f_\gamma t^\gamma:
 S\text{ well ordered and set-sized},\ f_\gamma\in\OO(U)\right\}.
\end{equation}
Coefficients that vanish identically are removed from the support.
For $F\ne0$, let $\nu(F)=\min\supp(F)$, with leading coefficient
$f_{\nu(F)}$. Products are Hahn products in the coefficient ring.
Because $U$ is connected, $\OO(U)$ is an integral domain and
$\nu(FG)=\nu(F)+\nu(G)$.

Define evaluation on $U^\mu$ by
\begin{equation}\label{eq:Heval}
 \ev(F)(a+\epsilon)=
 \sum_{\gamma\in S}\sum_{n\geq0}
    \frac{f_\gamma^{(n)}(a)}{n!}\epsilon^n t^\gamma.
\end{equation}
We generally write $F(z)$ for this evaluated function.

\begin{theorem}[Faithful coherent evaluation]\label{thm:evaluation}
The double family in \eqref{eq:Heval} is summable. Evaluation is an injective
algebra homomorphism into functions on $U^\mu$. For all $z\in U^\mu$,
\begin{equation}\label{eq:evalval}
 v(F(z))\geq\nu(F),\qquad
 \st\bigl(t^{-\nu(F)}F(a+\epsilon)\bigr)=f_{\nu(F)}(a).
\end{equation}
\end{theorem}
\begin{proof}
Let $E=\supp(\epsilon)$. Supports in the double sum lie in $S+E^*$.
The support lemma proves well-ordering and finite occurrence of each
exponent, including finiteness over Taylor orders $n$. Products can
therefore be regrouped. At an ordinary point $a$, the evaluated normal
form is $\sum f_\gamma(a)t^\gamma$. If evaluation is identically zero,
uniqueness of normal forms makes every $f_\gamma(a)$ zero for every $a$.
Thus every coefficient function vanishes. Finally, all terms arising
from $n>0$ have strictly increased exponent; the coefficient at the least
possible exponent is the leading coefficient evaluated at $a$.
\end{proof}
The hypothesis of a common ordinary domain is important. An arbitrary
family of holomorphic germs whose radii shrink to zero need not be an
element of $\HH(U)$ for any neighborhood $U$.

\subsection{Differentiation and Taylor's theorem}
Define the section derivative by
\[
 F'=\sum_{\gamma\in S}f_\gamma't^\gamma.
\]
\begin{theorem}[Full surreal Taylor theorem]\label{thm:taylor}
Let $F\in\HH(U)$ and $b=a+\epsilon\in U^\mu$. For every $h\in\m$,
\begin{equation}\label{eq:TaylorH}
 F(b+h)=\sum_{n\geq0}\frac{F^{(n)}(b)}{n!}h^n.
\end{equation}
The sum is Hahn summable. If $F\ne0$, then for every finite $N\geq0$
\begin{equation}\label{eq:remainder}
 F(b+h)-\sum_{n=0}^N\frac{F^{(n)}(b)}{n!}h^n
   =h^{N+1}Q_N(b,h),\qquad
 |Q_N(b,h)|<t^{\nu(F)-1}.
\end{equation}
Consequently the section derivative equals the full surreal derivative,
and all real and imaginary coordinate derivatives satisfy the
Cauchy--Riemann equations. Their real and imaginary parts are harmonic.
\end{theorem}
\begin{proof}
Use the ordinary formal Taylor identity for every $f_\gamma$ at $a$,
then substitute the two infinitesimals $\epsilon$ and $h$. The union of
their supports is positive and well ordered. Neumann's lemma and the
Minkowski-sum assertion justify both orders of summation. After factoring
$h^{N+1}$ from the tail, every remaining exponent is at least $\nu(F)$.
Apply \cref{lem:valuationbound}. For the derivative, take $N=1$ and choose
\[
 0<\delta<\min\{\omega^{-1},\eta\,t^{1-\nu(F)}\}.
\]
The difference quotient differs from $F'(b)$ by less than $\eta$ whenever
$0<|h|<\delta$. Taking real and imaginary increments gives the real-linear
first derivative. Repeating the calculation gives all second derivatives
and $u_{xx}+u_{yy}=v_{xx}+v_{yy}=0$.
\end{proof}
Here differentiation is with respect to the argument $z$. It keeps every
surreal scalar, including $t^\gamma$, constant. This is not a surreal-field
derivation of the kind studied in differential algebra or transseries.

\subsection{Coherence and analytic continuation}
\begin{proposition}[Coherent identity theorem]\label{prop:identity}
Suppose $U$ is connected and $F\in\HH(U)$.
If $F$ vanishes on one entire monad $a+\m$, then $F=0$.
More generally, if $F$ has zeros whose distinct standard parts accumulate
at an interior point of $U$ in the ordinary complex topology, then $F=0$.
\end{proposition}
\begin{proof}
If $F\ne0$ and $s=\nu(F)$, its leading coefficient $f_s$ is not
identically zero. Let $m$ be its finite vanishing order at $a$.
Then $F^{(m)}(a)$ has nonzero coefficient $f_s^{(m)}(a)$ at exponent $s$.
A function identically zero on a full neighborhood has all derivatives
zero, a contradiction. For the second statement, \eqref{eq:evalval}
implies that $f_s$ vanishes at all the indicated standard parts. The
classical identity theorem contradicts $f_s\ne0$.
\end{proof}
This does not assert existence of any nontrivial set-sized convergent
sequence in the full surreal topology. The accumulation condition is
explicitly on standard parts in $\C$.

Restriction of a section to a nonempty ordinary open subset is injective.
Compatible sections on a connected ordinary open cover glue uniquely:
on overlapping connected pieces their nonzero coefficient supports agree
by the ordinary identity theorem, and their coefficients glue holomorphically.
Connectivity propagates the common support through the overlap graph.
On disconnected domains, define sections component by component.
In this sense the construction gives a class-valued sheaf over the
\emph{ordinary complex plane}, not a claim of topological connectedness
of its surreal tubes.

\subsection{Units and infinitesimal functional calculus}
Write $\HH_{>0}(U)=\{F:\nu(F)>0\}\cup\{0\}$.
If $H\in\HH_{>0}(U)$, then
\begin{equation}\label{eq:smallcalc}
 (1+H)^{-1}=\sum_{n\geq0}(-H)^n,\quad
 \log(1+H)=\sum_{n\geq1}\frac{(-1)^{n+1}}nH^n,\quad
 e^H=\sum_{n\geq0}\frac{H^n}{n!}
\end{equation}
are well-defined sections. More generally, any ordinary formal power
series can be substituted into $H$. This is a direct application of the
positive-support lemma in the coefficient ring.

If $F=t^s(f_0+H)$ with $f_0$ nowhere zero on $U$ and $H\in\HH_{>0}(U)$,
then
\[
 F^{-1}=t^{-s}f_0^{-1}\bigl(1+H/f_0\bigr)^{-1}\in\HH(U).
\]
The condition is sufficient and is exactly what will be needed on
contour neighborhoods and for local preparation units.

\section{Residues before integration: an intrinsic algebraic theory}
\subsection{Formal meromorphic germs}
The ordinary formal Laurent field $\K((X))$ has a finite principal part
in its variable $X$. It is different from the Hahn expansion in the
surreal scale $t$. Define
\[
 \Res_{X=0}\left(\sum_{n\geq-N}a_nX^n\dd X\right)=a_{-1}.
\]
A formal meromorphic differential has a formal meromorphic primitive if
and only if its residue is zero: integrate every monomial except $X^{-1}$.
This is a purely algebraic assertion and does not assert evaluation of an
arbitrary Laurent series on an unspecified domain.

\begin{theorem}[Change of coordinates for residues]\label{thm:reschange}
If $X=\psi(Y)=c_1Y+c_2Y^2+\cdots$ with $c_1\ne0$, then
\[
 \Res_{Y=0} A(\psi(Y))\psi'(Y)\dd Y
   =\Res_{X=0} A(X)\dd X.
\]
\end{theorem}
\begin{proof}
For $A=X^n$ with $n\ne-1$, the substituted differential is the derivative
of $\psi^{n+1}/(n+1)$ and has zero residue. For $n=-1$, write
$\psi=c_1Y(1+B)$, where $B\in Y\K[[Y]]$. Then
\[
 \frac{\psi'}\psi=\frac1Y+\frac{B'}{1+B}
 =\frac1Y+\bigl(\log(1+B)\bigr)',
\]
whose residue is one. In the general Laurent series only finitely many
terms can contribute to any prescribed coefficient after substitution,
so the monomial argument applies.
\end{proof}
In particular, if $A=X^mU$ with $U(0)\ne0$, then
\begin{equation}\label{eq:localarg}
 \Res_{X=0}\frac{A'}A\dd X=m.
\end{equation}
This is the local argument principle, with an ordinary integer on its
right-hand side.

\subsection{The rational residue theorem on the surcomplex sphere}
Let $\mathbb P^1(\K)=\K\cup\{\infty\}$. For a rational differential
$R(z)\dd z$, define its residue at infinity using $w=1/z$:
\[
 \Res_\infty R(z)\dd z
 =\Res_{w=0}-R(1/w)w^{-2}\dd w.
\]
\begin{theorem}[Global rational residue theorem]\label{thm:rational-residue}
A rational differential over $\K$ has finitely many poles and
\[
 \sum_{b\in\mathbb P^1(\K)}\Res_b R(z)\dd z=0.
\]
It has a rational primitive if and only if all its residues vanish.
\end{theorem}
\begin{proof}
Algebraic closedness gives a finite partial-fraction expansion
\[
 R(z)=P(z)+\sum_{b}\sum_{j=1}^{m_b}\frac{c_{b,j}}{(z-b)^j}.
\]
The finite residues are $c_{b,1}$. The coefficient of $z^{-1}$ in the
expansion at infinity is $\sum_b c_{b,1}$, so the residue there is its
negative. Polynomial terms and terms with $j\geq2$ have rational
primitives. A derivative of a rational function has zero residue at every
point, proving the converse and necessity.
\end{proof}
No contour, completeness, or topological connectedness is involved.
This is a genuinely global theorem over the entire surcomplex field.

\section{Coefficientwise contours, Cauchy formulas, and primitives}
\subsection{What a Hahn contour integral means}
Let $\gamma$ be an ordinary piecewise $C^1$ complex path in $U$ and
$F=\sum f_s t^s\in\HH(U)$. Define
\begin{equation}\label{eq:Hintegral}
 \hint_\gamma F(\zeta)\dd\zeta
  :=\sum_s\left(\int_\gamma f_s(\zeta)\dd\zeta\right)t^s.
\end{equation}
The support can only shrink, so this is a well-defined surcomplex number.
The superscript $H$ is retained to emphasize that the operation is
\emph{coefficientwise}. It is not a limit of Riemann sums in the full
surreal topology. Linearity over $\K$ follows from the finite-decomposition
property of Hahn products.

The same construction works when the coefficient functions are merely
continuous along the path, or meromorphic on a common domain with no pole
on the path. It commutes with a summable family whose coefficientwise
integrals are defined, because each exponent involves only finitely many
members of that family.

\begin{theorem}[Cauchy's theorem]
If $\gamma$ is null-homologous in $U$ and $F\in\HH(U)$, then
$\hoint_\gamma F(\zeta)\dd\zeta=0$.
\end{theorem}
\begin{proof}
Apply the ordinary Cauchy theorem to every coefficient in
\eqref{eq:Hintegral}.
\end{proof}
This proof states exactly what is inherited from complex analysis. Its
use for points not in $\C$ requires an additional summability argument.

\subsection{Cauchy's formula at a surcomplex point}
\begin{theorem}[Surcomplex Cauchy formula]\label{thm:cauchy}
Let $\gamma=\partial\Omega$ be a positively oriented ordinary piecewise
$C^1$ Jordan curve, with $\overline\Omega\subset U$, and let
$z=a+\epsilon$ with $a\in\Omega$ and $\epsilon\in\m$. For every finite
$k\geq0$,
\begin{equation}\label{eq:CauchyH}
 \frac{F^{(k)}(z)}{k!}
 =\frac1{2\pi i}\hoint_\gamma
      \frac{F(\zeta)}{(\zeta-z)^{k+1}}\dd\zeta.
\end{equation}
On the right, the kernel is its Hahn expansion on an ordinary
neighborhood of $\gamma$.
\end{theorem}
\begin{proof}
Since $\zeta-a$ is nonzero near $\gamma$,
\[
 (\zeta-a-\epsilon)^{-k-1}
  =\sum_{n\geq0}\binom{n+k}{k}
       \frac{\epsilon^n}{(\zeta-a)^{n+k+1}}.
\]
The positive support of $\epsilon$ makes this a summable Hahn series with
holomorphic coefficients near $\gamma$. Multiply by $F$ and integrate
coefficientwise. The ordinary Cauchy derivative formula gives
\[
 \frac1{2\pi i}\hoint_\gamma
       \frac{F(\zeta)}{(\zeta-z)^{k+1}}\dd\zeta
 =\sum_{n\geq0}\frac{F^{(n+k)}(a)}{k!\,n!}\epsilon^n.
\]
The right side is $F^{(k)}(a+\epsilon)/k!$ by the Taylor theorem.
Every interchange is justified by summability, not by an unproved
convergence theorem.
\end{proof}

\subsection{Morera's criterion}
\begin{theorem}[Coefficientwise Morera theorem]\label{thm:morera}
Let $A=\sum_{s\in S}a_s t^s$ have continuous complex coefficient functions
on $U$, with a fixed well-ordered set support. If
\[
 \hoint_{\partial T} A(\zeta)\dd\zeta=0
\]
for every ordinary triangle whose closed interior lies in $U$, then each
$a_s$ is holomorphic. Consequently $A$ has a unique Hahn-coherent extension
to $U^\mu$.
\end{theorem}
\begin{proof}
Uniqueness of the normal form makes each coefficient integral zero.
Classical Morera applies to every continuous $a_s$. Evaluation then gives
an extension by \cref{thm:evaluation}, and faithfulness proves uniqueness
within the prescribed coefficient structure.
\end{proof}
Before the conclusion, the object is a Hahn-valued function on the ordinary
domain. No Taylor extension of a merely continuous coefficient is presumed.

\subsection{Primitives and their precise uniqueness}
\begin{theorem}[Coherent primitives]\label{thm:primitives}
If $U$ is simply connected, every $F\in\HH(U)$ has a primitive in
$\HH(U)$. Two coherent primitives differ by a constant in $\K$.
For a general connected $U$, a coherent primitive exists exactly when
all ordinary closed-path Hahn integrals vanish.
\end{theorem}
\begin{proof}
Fix $a_0\in U$. For every coefficient choose its ordinary primitive
$g_s$ with $g_s(a_0)=0$, using path independence. Then $G=\sum g_s t^s$
has support contained in $\supp(F)$ and satisfies $G'=F$.
If $H'=0$, uniqueness of normal forms gives $h_s'=0$ for each coefficient.
Connectedness of $U$ makes all $h_s$ constant, so $H$ is a surcomplex
constant. The period criterion follows by applying the same argument
coefficientwise in both directions.
\end{proof}
An arbitrary differentially holomorphic primitive may still be modified
by the locally constant functions of \cref{ex:locallyconstant}. Uniqueness
here belongs to the coherent category.

\section{Hahn--Weierstrass preparation}\label{sec:preparation}
\subsection{A division operator on a fixed disk}
Let $D$ be an ordinary disk centered at zero and $m\geq1$. For
$g\in\OO(D)$, define
\[
 T_mg=\sum_{j=0}^{m-1}\frac{g^{(j)}(0)}{j!}w^j,
 \qquad R_mg=\frac{g-T_mg}{w^m}.
\]
The apparent singularity of $R_mg$ at zero is removable. Thus
\[
 g=T_mg+w^mR_mg
\]
is a decomposition into a polynomial of degree less than $m$ and a
multiple of $w^m$. Apply these operators coefficientwise to Hahn sections;
they never lower the minimum exponent.

\begin{theorem}[Infinitesimal Hahn preparation]\label{thm:preparation}
Suppose
\[
 A(w)=w^m+H(w),\qquad H\in\HH_{>0}(D).
\]
There are unique objects
\[
 P(w)=w^m+\sum_{j=0}^{m-1}p_jw^j,
 \quad p_j\in\m,\qquad Q\in\HH_{>0}(D)
\]
such that
\begin{equation}\label{eq:prep}
 A(w)=P(w)\bigl(1+Q(w)\bigr).
\end{equation}
They are constructed using only coefficientwise Taylor division, products,
and summable families supported in the semigroup generated by
$\supp(H)$.
\end{theorem}
\begin{proof}
Write $P=w^m+p$. The desired identity is equivalent to
\[
 H=p+w^mQ+pQ.
\]
Applying $T_m$ and $R_m$ yields the simultaneous equations
\begin{equation}\label{eq:fixedprep}
 p=T_m(H-pQ),\qquad Q=R_m(H-pQ).
\end{equation}
Give each occurrence of $H$ formal degree one. Define
\[
 p^{[1]}=T_mH,\qquad Q^{[1]}=R_mH,
\]
and for $n\geq2$ define
\begin{align}\label{eq:preprecursion}
 p^{[n]}&=-T_m\sum_{j=1}^{n-1}p^{[j]}Q^{[n-j]},\\
 Q^{[n]}&=-R_m\sum_{j=1}^{n-1}p^{[j]}Q^{[n-j]}.
\end{align}
If $E=\supp(H)$, every degree-$n$ term is supported on sums of $n$
elements of $E$. For a fixed such word only finitely many binary
parenthesizations and operator applications arise. By Neumann's lemma,
\[
 p=\sum_{n\geq1}p^{[n]},\qquad Q=\sum_{n\geq1}Q^{[n]}
\]
are summable Hahn series. Their coefficients are holomorphic on the same
disk $D$ throughout the construction. The polynomial $p$ has degree less
than $m$. All exponents are positive. Summing the recursion proves
\eqref{eq:fixedprep} and therefore \eqref{eq:prep}.

For uniqueness, take two solutions and let $d$ be the minimum of the
valuations of their nonzero differences $p-p'$ and $Q-Q'$. Their
difference equations are obtained by applying $T_m,R_m$ to
\[
 -pQ+p'Q'=-(p-p')Q-p'(Q-Q').
\]
Because $p'$ and $Q$ have strictly positive valuation, the right side has
valuation strictly greater than $d$ whenever it is nonzero. The operators
cannot lower valuation. This contradicts the definition of $d$. If the
right side is zero, both differences are zero immediately.
\end{proof}
This is a formal-contraction proof using positive support. It is not a
Banach fixed-point proof and does not claim that its successive partial
sums converge in the full surreal topology.

\subsection{Preparation at a zero of a leading coefficient}
\begin{corollary}[Local preparation for a general section]\label{cor:generalprep}
Let $F\in\HH(U)$ be nonzero, $s=\nu(F)$, and suppose its leading
coefficient $f_s$ has a zero of order $m$ at $a\in U$. On a sufficiently
small ordinary disk $a+D\subset U$,
\begin{equation}\label{eq:generalprep}
 F(a+w)=t^s\left(w^m+\sum_{j=0}^{m-1}p_jw^j\right)V(w),
\end{equation}
where every $p_j$ is infinitesimal and $V\in\HH(D)$ has leading
coefficient $u_0(w)=f_s(a+w)/w^m$, which is nowhere zero on $D$.
In particular, $V$ is an invertible section and never vanishes on $D^\mu$.
\end{corollary}
\begin{proof}
Shrink $D$ so that $u_0$ has no zeros. Divide $t^{-s}F(a+w)$ by $u_0(w)$
to obtain $w^m+H$ with positive support. Apply
\cref{thm:preparation} and set $V=u_0(1+Q)$. Its inverse is supplied by
\eqref{eq:smallcalc}, and evaluation preserves the inverse identity.
\end{proof}
No uniform complex bound on the infinitely many perturbation coefficients
is required. The common holomorphic domain and support condition replace
that kind of bound.

\section{Zero clusters, Rouch\'e stability, and the argument principle}
\subsection{Exact counting inside a monad}
\begin{lemma}\label{lem:smallpolyroots}
All roots in $\K$ of a monic polynomial
\[
 P(w)=w^m+p_{m-1}w^{m-1}+\cdots+p_0,
 \qquad p_j\in\m,
\]
are infinitesimal.
\end{lemma}
\begin{proof}
If $\xi$ is not infinitesimal, then $\xi^{-1}$ is finite. Dividing
$P(\xi)=0$ by $\xi^m$ gives
$1+\sum_{j<m}p_j\xi^{j-m}=0$. Every summand after the one is
infinitesimal, an impossibility after taking standard parts.
\end{proof}

\begin{theorem}[Exact zero-cluster theorem]\label{thm:zerocluster}
Let $F\in\HH(U)$ be nonzero, with leading coefficient $f_s$.
For every $a\in U$, the number of zeros of $F$ in $a+\m$, counted
with their analytic multiplicities, is exactly
\[
 \ord_a(f_s).
\]
In particular, a nonzero leading value excludes all zeros in its monad;
a simple leading zero has exactly one surcomplex lift; and a leading
zero of order $m$ splits into exactly $m$ roots, with repetitions allowed.
\end{theorem}
\begin{proof}
If $f_s(a)\ne0$, \eqref{eq:evalval} excludes zeros. Otherwise apply
\cref{cor:generalprep}. The polynomial factor has exactly $m$ roots in
$\K$, counted algebraically, because $\K$ is algebraically closed.
By \cref{lem:smallpolyroots}, all these roots lie in $\m$. The section
$V$ does not vanish on the monad, so there are no other zeros.
The Taylor theorem and the product rule show that the order of vanishing
of $P(w)V(w)$ at each root is precisely its polynomial multiplicity.
\end{proof}
This gives an actual finite zero count in a proper-class neighborhood.
The count is not a reduction modulo infinitesimals: root positions may
contain arbitrarily deep surreal terms, while the count itself is exact.

\begin{corollary}[Finite-domain zero count]\label{cor:domaincount}
Suppose $\Omega$ is an ordinary bounded Jordan domain with
$\overline\Omega\subset U$, and $f_s$ has no zero on $\partial\Omega$.
Then $F$ has finitely many zeros in $\Omega^\mu$, and their total
multiplicity equals the ordinary zero count of $f_s$ in $\Omega$.
\end{corollary}
\begin{proof}
A zero of $F$ can occur only above a zero of $f_s$. There are finitely
many such standard zeros in $\Omega$, since $f_s$ is holomorphic near
its compact closure and nonzero on its boundary. Sum
\cref{thm:zerocluster} over these points.
\end{proof}

\subsection{Two useful Rouch\'e statements}
\begin{theorem}[Hahn Rouch\'e stability]\label{thm:rouche}
Under the hypotheses of \cref{cor:domaincount}, adding any
$G\in\HH(U)$ with $\nu(G)>\nu(F)$ leaves the total number of zeros in
$\Omega^\mu$ unchanged, with multiplicity.

More generally, suppose
\[
 F=t^s(f_0+F_+),\qquad G=t^s(g_0+G_+),\qquad
 F_+,G_+\in\HH_{>0}(U),
\]
and
$|g_0-f_0|<|f_0|$ on $\partial\Omega$. Then $F$ and $G$ have the same
zero count in $\Omega^\mu$.
\end{theorem}
\begin{proof}
In the first assertion the leading coefficient does not change. Apply
\cref{cor:domaincount}. In the second, ordinary Rouch\'e gives equal
zero counts for $f_0,g_0$ and excludes boundary zeros of both; the same
corollary lifts those counts.
\end{proof}
The strict inequality in the second assertion is an ordinary complex
margin. It is not silently weakened to a pointwise non-strict inequality
between standard parts.

\subsection{Hahn residues of meromorphic coefficients}
Let $\MM(U)$ consist of set-supported Hahn series with coefficients
meromorphic on the connected ordinary domain $U$. Since the coefficient
ring is now a field, nonzero elements admit inverses by leading-term
factorization and \eqref{eq:smallcalc}. Define
\[
 \Res_a^H A(\zeta)\dd\zeta
   =\sum_s\bigl(\Res_a a_s(\zeta)\dd\zeta\bigr)t^s.
\]
This is a \emph{cluster residue}: it is attached to an ordinary center,
not necessarily to a single surcomplex pole.

\begin{theorem}[Hahn residue theorem]\label{thm:Hresidue}
Suppose $A=\sum a_s t^s\in\MM(U)$, $\overline\Omega\subset U$, and no
coefficient $a_s$ has a pole on $\gamma=\partial\Omega$. Then
\begin{equation}\label{eq:Hresidue}
 \frac1{2\pi i}\hoint_\gamma A(\zeta)\dd\zeta
   =\sum_{a\in\Omega}\Res_a^H A(\zeta)\dd\zeta.
\end{equation}
Only a set of centers contributes, and the family on the right is
Hahn summable even when infinitely many centers contribute.
\end{theorem}
\begin{proof}
For each exponent, its coefficient has finitely many poles in
$\overline\Omega$, by meromorphicity on a neighborhood of that compact
set. Therefore only finitely many centers contribute to a fixed
exponent. The union of all residue supports is contained in the fixed
well-ordered support of $A$. This proves summability and permits summing
the ordinary residue theorem coefficientwise.
\end{proof}
For example, if $\epsilon\in\m$, the rational function
$1/(\zeta-\epsilon)$ has the Hahn expansion
$\sum_{n\geq0}\epsilon^n\zeta^{-n-1}$ away from the zero standard-part
monad. Its cluster residue at zero is one, correctly accounting for the
actual pole at $\epsilon$. This expansion is not asserted to evaluate
throughout the monad containing that pole.

\begin{theorem}[Exact argument principle]\label{thm:argument}
Under the hypotheses of \cref{cor:domaincount},
\begin{equation}\label{eq:argument}
 \frac1{2\pi i}\hoint_{\partial\Omega}\frac{F'(\zeta)}{F(\zeta)}\dd\zeta
 =\sum_{z\in\Omega^\mu}\ord_zF
 =\sum_{a\in\Omega}\ord_a f_s\in\N.
\end{equation}
The integral has no infinitesimal correction. More locally,
$\Res_a^H(F'/F)\dd\zeta=\ord_a f_s$.
\end{theorem}
\begin{proof}
Write $F=t^sf_s(1+H)$ in the meromorphic coefficient field, with
$\nu(H)>0$. Then
\[
 \frac{F'}F=\frac{f_s'}{f_s}+\bigl(\log(1+H)\bigr)'.
\]
Each coefficient of the second summand is the derivative of a
meromorphic function and therefore has zero residue everywhere and zero
integral on a closed contour avoiding its poles. The ordinary argument
principle gives the integral and local residue of the first summand.
The exact zero-cluster theorem identifies these integers with the
surcomplex multiplicities.
\end{proof}

\section{Local mapping theorems at arbitrary surcomplex points}
\subsection{A normalization lemma}
\begin{lemma}[Infinitesimal rescaling]\label{lem:normalize}
Let $F\in\HH(U)$ be nonconstant and $b\in U^\mu$. There is a least
$m\geq1$ for which $a_m=F^{(m)}(b)/m!\ne0$. For a sufficiently small
positive monomial $r=t^\rho$, the normalized function
\begin{equation}\label{eq:normalize}
 G(w)=\frac{F(b+rw)-F(b)}{a_mr^m}
\end{equation}
is the evaluation of a section of $\HH(\C)$ of the form
$w^m+H(w)$ with $\nu(H)>0$. The monomial $r$ may be chosen small enough
that $b+r\m$ lies in any specified full-open neighborhood of $b$.
\end{lemma}
\begin{proof}
First apply the coherent identity theorem to $F-F(b)$: not all of its
derivatives at $b$ can vanish. More explicitly, its first nonzero leading
coefficient has finite order at the standard part of $b$, so some
positive-order derivative is nonzero.

Write $b=a+\epsilon$, and choose a lower bound $\sigma$ for the exponents
in the common supports of the Taylor coefficients $F^{(n)}(b)/n!$.
One may take $\sigma=\nu(F)$ when $F\ne0$. Those supports are all
contained in the same well-ordered set
$\supp(F)+\supp(\epsilon)^*$. Choose
\[
 \rho>0,\qquad \rho>v(a_m)-\sigma.
\]
The Taylor expansion in \cref{thm:taylor}, with $rw$ infinitesimal for
finite $w$, gives
\[
 G(w)=w^m+\sum_{n>m}\frac{a_n}{a_m}r^{n-m}w^n.
\]
The supports of the tail lie in the sum of a fixed well-ordered set and
$\{\rho,2\rho,3\rho,\ldots\}$. Each exponent receives contributions from
only finitely many $n$, so its coefficient is a polynomial in $w$.
Its smallest possible exponent is at least
$\sigma-v(a_m)+\rho>0$. This proves the asserted entire Hahn
representation. Increasing $\rho$ makes $r$ smaller than any prescribed
positive radius, and hence puts $b+r\m$ inside the given neighborhood.
\end{proof}
The lemma is particularly useful when $F'(b)$ is nonzero but itself
infinitesimal. Looking only at the derivative of the leading complex
coefficient would miss that case.

\subsection{Open mapping and finite local degree}
\begin{theorem}[Local finite-degree mapping]\label{thm:open}
With the notation of \cref{lem:normalize},
\[
 F(b+r\m)=F(b)+a_mr^m\m.
\]
Every value on the right has exactly $m$ preimages in $b+r\m$, counted
with multiplicity. Consequently every nonconstant section in $\HH(U)$
is an open map for the full surreal topology.
\end{theorem}
\begin{proof}
For each $\eta\in\m$, the leading coefficient of $G(w)-\eta$ is $w^m$.
The zero-cluster theorem gives exactly $m$ solutions in $\m$.
Conversely, $G(w)$ is infinitesimal for $w\in\m$, by its positive-support
normalization. Undo the affine changes to obtain the equality and count.
For an arbitrary full-open neighborhood of $b$, choose $r$ as in the
last assertion of \cref{lem:normalize}. Its image contains the full-open
neighborhood $F(b)+a_mr^m\m$ of $F(b)$.
\end{proof}
This is a statement about actual open classes and actual surcomplex
preimages, not only about the reduction to $\C$.

\begin{corollary}[Strong local maximum principle]\label{cor:max}
The absolute value of a nonconstant Hahn-coherent holomorphic section
has no local maximum in the full surreal topology.
\end{corollary}
\begin{proof}
Every neighborhood of an image value contains values of strictly larger
absolute value. For a nonzero value $c$, use $(1+\eta)c$ with a sufficiently
small positive surreal $\eta$; for zero use any sufficiently small
nonzero value. This contradicts a local maximum because the map is open.
\end{proof}

\subsection{An inverse in the coherent category}
\begin{lemma}[Hahn reversion]\label{lem:reversion}
If $G(w)=w+H(w)$ with $H\in\HH_{>0}(\C)$, there is a unique
$K\in\HH_{>0}(\C)$ such that $J(y)=y+K(y)$ is a two-sided inverse
under infinitesimal Hahn composition. Its evaluated function is a
bijection $\fin\to\fin$. One explicit formula is
\begin{equation}\label{eq:reversion}
 K(y)=\sum_{n\geq1}\frac{(-1)^n}{n!}
       \left(\frac{\mathrm d}{\mathrm dy}\right)^{n-1}H(y)^n.
\end{equation}
\end{lemma}
\begin{proof}
Expand $H(y+K)$ as the formal Taylor series in the positive-support
section $K$. The equation
\[
 K=-H(y+K)
\]
determines the part of degree $n$ in occurrences of $H$ from the parts of
smaller degree. Derivatives of entire coefficient functions stay entire.
The support of a degree-$n$ term is a sum of $n$ elements of
$E=\supp(H)$. Neumann's lemma makes the recursively constructed terms
summable. A least-exponent comparison proves uniqueness, since changing
$K$ on the right multiplies its difference by a positive-support section.

To obtain the displayed formula, insert an auxiliary ordinary formal
parameter $q$ and solve $u=-qH(y+u)$. Formal Lagrange inversion, proved by
coefficient extraction in Appendix~\ref{app:lagrange}, gives
\[
 [q^n]u=\frac{(-1)^n}{n!}D_y^{n-1}H(y)^n.
\]
Summability permits setting $q=1$ in the graded sense, proving
\eqref{eq:reversion}.
For a finite target $y=c+\eta$, the zero-cluster theorem applied to
$G(w)-y$ gives precisely one solution in $c+\m$, and any finite solution
must have that standard part. Thus evaluation is bijective. The constructed
right inverse is that bijection's inverse and is also a left inverse;
faithful evaluation gives the corresponding section identity.
\end{proof}

\begin{theorem}[Surcomplex local inverse theorem]\label{thm:inverse}
If $F\in\HH(U)$ and $F'(b)\ne0$, then $F$ has a local inverse between
full-open neighborhoods of $b$ and $F(b)$. After affine rescaling, the
inverse is Hahn-coherent, and
\[
 (F^{-1})'(F(z))=\frac1{F'(z)}
\]
on a sufficiently small such neighborhood.
\end{theorem}
\begin{proof}
Use $m=1$ in \cref{lem:normalize}, then \cref{lem:reversion} and undo
the affine normalizations. The derivative formula follows from the chain
rule. A still smaller neighborhood makes the derivative nonzero
throughout, either by its first-order estimate or by the normalized
leading derivative one.
\end{proof}

\section{Cauchy estimates, Schwarz's lemma, and the limits of bounds}
\subsection{Coefficientwise Cauchy majorants}
Let the closed ordinary disk $\overline{D(a,r)}$ lie in $U$, where $r>0$
is real. For $F=\sum f_s t^s$, put
\[
 M_s=\max_{|\zeta-a|=r}|f_s(\zeta)|,
 \qquad B_r(F)=\sum_s M_s t^s\in\No^{\geq0}.
\]
These are ordinary compact maxima, followed by a Hahn sum.

\begin{proposition}[A Hahn Cauchy estimate]\label{prop:estimate}
For every finite $n\geq0$,
\[
 |F^{(n)}(a)|\leq \frac{n!}{r^n}B_r(F).
\]
\end{proposition}
\begin{proof}
Classical Cauchy estimates give
$|f_s^{(n)}(a)|\leq n!M_s/r^n$ coefficientwise. It remains to justify
that a coefficientwise absolute-value majorant bounds the surreal
absolute value. If $Z=\sum c_s t^s$ and $B=\sum |c_s|t^s$, then
\[
 B^2-|Z|^2
 =\sum_{s,u}\bigl(|c_s||c_u|-\Re(c_s\overline{c_u})\bigr)t^{s+u}
 \geq0,
\]
because all coefficients in this summable family are nonnegative real
numbers. Therefore $|Z|\leq B$. Increasing the nonnegative coefficients
preserves the order, and proves the estimate.
\end{proof}
The majorant $B_r(F)$ is not generally the least upper bound of $|F|$ on
a surreal circle. Its coefficientwise definition is part of the statement.

\subsection{A sharp Schwarz lemma for canonical lifts}
\begin{theorem}[Lifted Schwarz lemma]\label{thm:schwarz}
Let $f:\Disk\to\Disk$ be ordinary holomorphic with $f(0)=0$. Then
\[
 |\widetilde f(z)|\leq|z|\qquad(z\in\Disk^\mu).
\]
Equality at a nonzero $z$ occurs if and only if $f(w)=cw$ for a complex
constant $c$ with $|c|=1$.
\end{theorem}
\begin{proof}
Write $f(w)=wg(w)$ with $g\in\OO(\Disk)$. By ordinary Schwarz's lemma,
either $g$ is a constant of modulus one or $|g(a)|<1$ for every
$a\in\Disk$. In the second case,
$\st(\widetilde g(z))=g(\st z)$, so
$|\widetilde g(z)|<1$: an infinitesimal perturbation cannot remove a
strict ordinary real margin. Now use
$\widetilde f(z)=z\widetilde g(z)$. The constant case gives equality.
\end{proof}
The restriction to canonical lifts is necessary. For any positive
$\epsilon\in\m$, the coherent section $F(z)=(1+\epsilon)z$ maps
$\Disk^\mu$ into itself, fixes zero, and satisfies $|F(z)|>|z|$ for
nonzero $z$. Thus the same sharp inequality is false for all coherent
self-maps of the tube. The tube does not include points infinitesimally
inside the boundary whose standard part lies on that boundary.

\subsection{Why finite-tube boundedness is too weak for Liouville}
Every nonzero $F\in\HH(\C)$ satisfies
\begin{equation}\label{eq:automaticbound}
 |F(z)|<t^{\nu(F)-1}\qquad(z\in\C^\mu=\fin),
\end{equation}
by \eqref{eq:evalval} and \cref{lem:valuationbound}. Hence allowing an
arbitrary surreal bound makes boundedness on the finite tube automatic.
Even a real bound does not suffice: $\epsilon\widetilde{e^z}$ is
nonconstant and has absolute value less than one at every finite
surcomplex argument.

\begin{proposition}[Coefficientwise Liouville theorem]
If $F=\sum f_s t^s\in\HH(\C)$ and every $f_s$ is bounded on $\C$
by an ordinary real constant (depending on $s$), then $F$ is constant.
\end{proposition}
\begin{proof}
Classical Liouville makes each coefficient constant. Equivalently, apply
ordinary Cauchy estimates and let the ordinary radius tend to infinity
separately for each coefficient.
\end{proof}
There is no ordinary limiting process through radii cofinal in all of
$\No$. A genuinely global replacement is supplied by the next section.

\section{All-scale coherence forces polynomiality}\label{sec:rigidity}
\subsection{A global condition with an exact classification}
Fix the ordinary unit disk $\Disk$. For a class function $F:\K\to\K$,
consider the following condition:
\begin{equation}\label{eq:allscale}
 \text{for every }R\in\No^{>0},\quad
 w\longmapsto F(Rw)\text{ on }\Disk^\mu
 \text{ is the evaluation of some }G_R\in\HH(\Disk).
\end{equation}
There need not be one common support for all $R$; each scale may use its
own set-supported Hahn representation. Even with that freedom the
condition is highly restrictive.

\begin{theorem}[All-scale coherence rigidity]\label{thm:rigidity}
A function $F:\K\to\K$ satisfies \eqref{eq:allscale} if and only if
it is a polynomial in $z$ with coefficients in $\K$ and ordinary finite
degree.
\end{theorem}
\begin{proof}
A polynomial clearly satisfies the condition: a finite union of its
coefficient supports remains well ordered, and the coefficient functions
in $w$ are polynomials.

Conversely, the chart at $R=1$ supplies derivatives of every order at zero.
Put $a_n=F^{(n)}(0)/n!$. The sequence $(a_n)_{n\in\N}$ is set-sized.
Suppose infinitely many $a_n$ are nonzero. The surreal cut property gives
$A>0$ such that
\[
 |v(a_n)|<A\qquad(a_n\ne0).
\]
Choose $D>2A$ and the positive surreal scale $R=\omega^D$, so
$v(R)=-D$. For increasing indices $n<n'$ with nonzero coefficients,
\[
 v(a_{n'}R^{n'})-v(a_nR^n)
   =v(a_{n'})-v(a_n)-(n'-n)D<2A-D<0.
\]
Thus the valuations of the nonzero $a_nR^n$ form an infinite strictly
decreasing sequence.

On the other hand, write the coherent representative at that scale as
$G_R(w)=\sum_{s\in S_R}g_s(w)t^s$. Differentiation and the finite chain
rule give
\[
 \frac{G_R^{(n)}(0)}{n!}=a_nR^n.
\]
Every one of these values has normal-form support contained in the same
well-ordered set $S_R$. In particular its valuation, when nonzero, is
an element of $S_R$. The strictly decreasing sequence just constructed
contradicts well-ordering.

Therefore $a_n=0$ for all $n>N$ for some finite $N$. At any scale $R$,
uniqueness of normal forms in the derivative identities makes
$g_s^{(n)}(0)=0$ for every $s$ and every $n>N$. Ordinary holomorphy and
the ordinary identity theorem give
\[
 g_s(w)=\sum_{n=0}^N\frac{g_s^{(n)}(0)}{n!}w^n
 \quad(w\in\Disk).
\]
Interchanging this finite sum with the Hahn sum, then evaluating, yields
\[
 F(Rw)=\sum_{n=0}^N a_nR^nw^n\qquad(w\in\Disk^\mu).
\]
For any $z\in\K$, choose a positive $R>2|z|+1$. Then $z/R$ has standard
part in the ordinary unit disk, so the last identity gives
$F(z)=\sum_{n=0}^N a_nz^n$.
\end{proof}
The key point is not an Archimedean radius estimate. Every set of
coefficient valuations can be dominated by one surreal scale, and that
scale would force an infinite descending support if infinitely many
Taylor coefficients survived.

\subsection{Global Liouville and polynomial growth}
\begin{corollary}[All-scale Liouville theorem]\label{cor:globalLiouville}
If $F$ satisfies \eqref{eq:allscale} and $|F(z)|\leq M$ for all
$z\in\K$, for some $M\in\No$, then $F$ is constant.
More generally, if
\[
 |F(z)|\leq C(1+|z|)^N\quad(z\in\K)
\]
with $C>0$ surreal and $N$ an ordinary nonnegative integer, then
$\deg F\leq N$.
\end{corollary}
\begin{proof}
By the theorem, $F$ is a polynomial. If its degree $d$ is greater than
$N$, choose a sufficiently large positive real-coordinate surreal $x$
that the leading term dominates all lower terms and
$|a_d|x^{d-N}>2^{N+1}C$. Then
$|F(x)|\geq |a_d|x^d/2>C(1+x)^N$, a contradiction. The case $N=0$
gives the first assertion after enlarging a nonnegative bound if needed.
\end{proof}
This condition should not be called the unique definition of entire
surcomplex analyticity. It excludes the canonical exponential, whose
Taylor coefficients at zero are all nonzero. It instead precisely
classifies what is possible when one demands the particular Hahn
coherence of this article on \emph{every} scale.

\section{The global surcomplex exponential and logarithms}\label{sec:exponential}
\subsection{A normal-form description of the canonical exponential}
For a real surreal $y$, split its normal form uniquely as
\[
 y=y_{\mathrm{inf}}+y_0+\eta,
\]
where $y_{\mathrm{inf}}$ contains only infinite monomials (negative
$t$-exponents), $y_0\in\R$, and $\eta$ is infinitesimal. The purely
infinite parts form an additive real vector space, denoted $\mathcal J$.
The three projections are additive. They are not standard-part maps on
all of $\No$: standard part is reserved for finite numbers.

We use the Gonshor real surreal exponential $\exp_{\No}$ and its inverse
$\log_{\No}$. In addition to the exponential group law, its infinitesimal
Taylor identity is
\[
 \exp_{\No}(u)=\sum_{n\geq0}\frac{u^n}{n!}\qquad(u\in\m\cap\No).
\]
These are established properties of the surreal exponential
\cite{Gonshor,vdDE}.

Define
\begin{equation}\label{eq:canonicalexp}
 \Exp_0(x+iy)
  =\exp_{\No}(x)\,e^{iy_0}
       \sum_{n\geq0}\frac{(i\eta)^n}{n!}.
\end{equation}
The ordinary complex exponential is used only at the real number $y_0$;
the last factor is a Hahn sum. Purely infinite imaginary parts do not
contribute to this canonical phase. This is the normal-form version of
the canonical Ehrlich--Kaplan exponential \cite[Theorem 11.1]{EK}.
The corresponding omnific integers are
\[
 \Oz=\mathcal J+\Z.
\]

\begin{theorem}[Global exponential properties]\label{thm:exp}
The map $\Exp_0:\K\to\K^\times$ is a surjective additive-to-multiplicative
group homomorphism. It extends the ordinary complex exponential and
$\exp_{\No}$, and satisfies
\begin{equation}\label{eq:exp-props}
 |\Exp_0(x+iy)|=\exp_{\No}(x),\qquad
 \ker\Exp_0=2\pi i\Oz.
\end{equation}
It is differentially holomorphic everywhere, with
\[
 \Exp_0'=\Exp_0,
 \qquad
 \Exp_0(b+h)=\Exp_0(b)\sum_{n\geq0}\frac{h^n}{n!}
 \quad(h\in\m).
\]
\end{theorem}
\begin{proof}
The additivity of the normal-form projections, the real exponential group
law, and the formal identity $e^{u+v}=e^ue^v$ for infinitesimals prove
the group law. Conjugating the final factor in \eqref{eq:canonicalexp}
replaces $\eta$ by $-\eta$, so that factor has absolute value one.
This proves the modulus identity.

If the value is one, its modulus forces $x=0$. Its standard phase is
$e^{iy_0}=1$, so $y_0\in2\pi\Z$. The first nonzero term of
$e^{i\eta}-1=i\eta+\cdots$ cannot vanish unless $\eta=0$.
Thus the imaginary kernel is $\mathcal J+2\pi\Z=2\pi\Oz$.

For surjectivity, let $q\ne0$, put $\rho=|q|$, and set $u=q/\rho$.
The number $u$ is finite, has modulus one, and has standard part
$c\in\C$ of modulus one. Write $u=c(1+\delta)$ with $\delta\in\m$.
The infinitesimal logarithm
$\ell=\sum_{n\geq1}(-1)^{n+1}\delta^n/n$ is defined. Since
$(1+\delta)(1+\bar\delta)=1$, the formal logarithm identity implies
$\ell+\bar\ell=0$. Therefore $\ell=i\eta$ for a real infinitesimal
$\eta$. Choose an ordinary real $\theta$ with $c=e^{i\theta}$.
Then \eqref{eq:canonicalexp} sends
$\log_{\No}\rho+i(\theta+\eta)$ to $q$.

For an infinitesimal complex $h$, its real and imaginary Taylor
exponentials multiply to the formal series $\sum h^n/n!$.
The group law gives the displayed local expansion at $b$. Factoring
$h$ from its linear remainder bounds the difference quotient error by
$|\Exp_0(b)|\omega|h|$, which tends to zero in the full
$\eta$--$\delta$ sense. This proves the derivative assertion.
\end{proof}
Although global, this exponential does not have a Hahn-summable Maclaurin
series at every argument. For an infinite positive real argument it is
defined by $\exp_{\No}$, not by declaring $\sum z^n/n!$ summable.

\subsection{Local logarithm branches}
\begin{theorem}[Infinitesimal logarithmic charts]\label{thm:log}
Let $b\in\K^\times$ and choose $\lambda$ with $\Exp_0(\lambda)=b$.
For $w\in b(1+\m)$ define
\[
 \Log_\lambda(w)=\lambda+
   \sum_{n\geq1}\frac{(-1)^{n+1}}n\left(\frac{w-b}{b}\right)^n.
\]
This is a bijection $b(1+\m)\to\lambda+\m$, inverse to the exponential
on that monad, and $\Log_\lambda'(w)=1/w$.
Any two choices of $\lambda$ differ by an element of $2\pi i\Oz$.
\end{theorem}
\begin{proof}
The formal identities $\log(e^h)=h$ and $e^{\log(1+h)}=1+h$ evaluate
at infinitesimals. The derivative is the evaluated derivative of the
formal logarithm, followed by the affine chain rule. The last assertion
is the kernel calculation in \cref{thm:exp}.
\end{proof}
This describes the branch choices without imposing an unmentioned
global principal argument or a path-connected covering-space structure.

\subsection{The differential equation does not select a global phase}
For $y\in\No$, let $L_\omega(y)\in\R$ denote the coefficient of the
normal-form monomial $\omega$ (equivalently, exponent $-1$ in $t$).
This is an additive real-linear map, and it is zero on every finite $y$.
For $c\in\R$, define
\begin{equation}\label{eq:twistedexp}
 \Exp_c(x+iy)=\Exp_0(x+iy)\,
           \exp_{\C}\bigl(icL_\omega(y)\bigr).
\end{equation}

\begin{theorem}[Explicit differential nonuniqueness]\label{thm:twists}
Every $\Exp_c$ is a surjective exponential homomorphism extending
$\exp_{\No}$ and the ordinary complex exponential. It agrees with
$\Exp_0$ on all finite surcomplex arguments, satisfies
\[
 |\Exp_c(x+iy)|=\exp_{\No}(x),\qquad \Exp_c'=\Exp_c,
\]
and is locally Hahn-coherent in finite translated charts.
For $c\notin2\pi\Z$ it differs from $\Exp_0$, since
$\Exp_c(i\omega)=e^{ic}$ whereas $\Exp_0(i\omega)=1$.
\end{theorem}
\begin{proof}
The correction factor is a character of the additive imaginary group
with values in the ordinary unit circle. It is one on finite imaginary
parts, so all extension claims and the modulus identity follow.
Surjectivity follows from the preimages constructed in \cref{thm:exp},
whose imaginary parts are finite and therefore unchanged by the twist.
For an infinitesimal increment $h$, the correction factor does not change:
$L_\omega(\Im h)=0$. Hence the same local exponential Taylor identity
and derivative proof apply. On each translated finite tube $b+\fin$,
\[
 \Exp_c(b+w)=\Exp_c(b)\widetilde{e^w},\qquad w\in\fin,
\]
a product of a surcomplex constant and a canonical lift, so it is
coherent in that normalized chart. The special value is immediate.
\end{proof}
In particular, $E'=E$, $E(0)=1$, agreement with the real surreal
exponential, and the correct modulus law do not determine a unique
global surcomplex exponential. Requiring the phase to annihilate the
purely infinite imaginary summand does select \eqref{eq:canonicalexp}.
The twists do not assert a solution to the additional integer-part
robustness questions formulated by Ehrlich and Kaplan; those questions
contain structural hypotheses beyond these differential conditions.

\begin{corollary}
For $c\notin2\pi\Z$, the quotient $\Exp_c/\Exp_0$ is a nonconstant
unit-modulus function on all of $\K$, differentially holomorphic with
all positive-order derivatives zero.
\end{corollary}
This supplies a more structured version of the locally constant
counterexample. Local analytic Taylor formulas alone do not impose global
analytic continuation between infinitely separated parts of $\K$.

\section{A Picard-type theorem at an explicit essential singularity}
Define $P(z)=\Exp_0(1/z)$ for $z\ne0$. It is differentially holomorphic
on $\K^\times$, with $P'(z)=-z^{-2}P(z)$. It never takes the value zero.

\begin{theorem}[Every nonzero value at every surreal scale]\label{thm:picard}
For every $w\in\K^\times$ and every positive surreal radius $\delta$,
there is a proper class of distinct $z$ satisfying
\[
 0<|z|<\delta,\qquad \Exp_0(1/z)=w.
\]
The function has neither a removable singularity nor a pole at zero in
the full surreal topology.
\end{theorem}
\begin{proof}
Choose a logarithm $\lambda$ of $w$. For any $q\in\Oz$ with
$\lambda+2\pi iq\ne0$, put
\[
 z_q=\frac1{\lambda+2\pi iq}.
\]
These are distinct and satisfy $P(z_q)=w$. Omnific integers are cofinal
in the positive surreals; indeed they contain all ordinals. For every
sufficiently large positive $q$,
\[
 |\lambda+2\pi iq|\geq2\pi q-|\lambda|>\delta^{-1}.
\]
A proper-class tail of such $q$ gives the required solutions.
Taking $w=1$ and $w=2$ shows that no limit at zero exists: every punctured
neighborhood contains both values. Taking $w=1$ also excludes a pole,
which would require the absolute value to exceed every fixed bound near
zero.
\end{proof}
This is an exact value-distribution theorem for a specified function,
not a general Picard theorem for all differentially holomorphic functions.
A general assertion of that breadth would be incompatible with the
locally constant examples. Nor does the phrase ``essential singularity''
here assert the existence of a global Laurent classification in the
unrestricted differential category: it records the explicitly proved
failure of removability and of a pole.

On ordinary annuli avoiding zero, $P$ agrees with the canonical lift of
$e^{1/z}$. Its coefficientwise contour integral on an ordinary positive
circle is therefore $2\pi i$, because the ordinary Laurent coefficient
of $z^{-1}$ is one. This is consistent with, but not defined by, its
behavior on infinitesimal monads near the singularity.

\section{Changes of scale and several variables}
\subsection{Affine charts at infinite and infinitesimal locations}
For $b\in\K$ and $r\in\K^\times$, a scaled tube is
\[
 b+rU^\mu=\{b+rw:w\in U^\mu\}.
\]
A section $F\in\HH(U)$ defines a function there by $z\mapsto F((z-b)/r)$.
Derivatives acquire factors $r^{-n}$. For differential forms, the
coefficientwise contour convention in that chart is
\[
 \hint_{b+r\gamma} A(z)\dd z
 :=\hint_\gamma A(b+rw)\,r\dd w.
\]
This notation denotes a chart-defined Hahn contour, not a continuous
ordinary path into the full surreal topology. Residues of differentials
are invariant under the change, by \cref{thm:reschange}.

There is a useful explicit compatibility criterion. Suppose the normalized
coordinate change is $z=\alpha+\beta w$, with $\alpha,\beta$ finite.
Write $\alpha=\alpha_0+\delta_\alpha$ and
$\beta=\beta_0+\delta_\beta$, with ordinary complex standard parts and
infinitesimal remainders. If $\alpha_0+\beta_0V\subset U$, then
\begin{equation}\label{eq:chartchange}
 F(\alpha+\beta w)=
 \sum_{s,n}\frac{f_s^{(n)}(\alpha_0+\beta_0w)}{n!}
       (\delta_\alpha+\delta_\beta w)^n t^s
\end{equation}
defines a section of $\HH(V)$. The case $\beta_0=0$ is allowed, provided
$\alpha_0\in U$. All support increases come from the fixed positive
supports of $\delta_\alpha,\delta_\beta$, so the proof is again Neumann's
lemma. An infinite normalized coefficient $\beta$ does not satisfy this
criterion. That distinction is precisely why an exponential can be
coherent locally without satisfying all-scale coherence.

\subsection{A multivariable implicit-function theorem}
For $U\subset\C^p$, define $\HH(U)$ by holomorphic coefficient
functions of $p$ complex variables. Evaluation at $a+\epsilon$ uses
\[
 F(a+\epsilon)=\sum_s\sum_{\alpha\in\N^p}
       \frac{\partial^\alpha f_s(a)}{\alpha!}\epsilon^\alpha t^s.
\]
Only finitely many infinitesimal supports are involved, so the same support
lemma proves summability, Taylor formulas, and the complex chain rule.

\begin{theorem}[Hahn implicit-function theorem]\label{thm:implicit}
Let $F=(F_1,\ldots,F_q)$ be a vector of Hahn sections on an ordinary
neighborhood of $(a,b)\in\C^p\times\C^q$. Suppose all exponents are
nonnegative, the exponent-zero vector is $f_0$, and
\[
 f_0(a,b)=0,\qquad \det\partial_y f_0(a,b)\ne0.
\]
Let $g_0$ be the ordinary implicit solution near $a$. On a sufficiently
small ordinary neighborhood $U'$ of $a$, there is a unique solution
\[
 g=g_0+h,\qquad h\in\HH_{>0}(U')^q,
 \qquad F(x,g(x))=0.
\]
Evaluation gives the unique solution in the corresponding $g_0$-monad
for every $x\in(U')^\mu$.
\end{theorem}
\begin{proof}
Shrink $U'$ so that $A(x)=\partial_y f_0(x,g_0(x))$ is invertible
holomorphically. Taylor expansion in $h$ gives
\[
 F(x,g_0+h)=A(x)h+B(x)+N(x,h),
\]
where $B$ has positive support, and every term of $N$ either is at least
quadratic in $h$ or contains a positive-support coefficient and at least
one occurrence of $h$. Solve
\[
 h=-A^{-1}(B+N(x,h))
\]
recursively in the number of positive-support coefficients, exactly as
in preparation and reversion. Each recursion coefficient is holomorphic
on the same smaller domain. Neumann's lemma makes the tree expansion
summable, and comparison of the least exponent gives uniqueness.

For pointwise uniqueness, take two solutions with the same prescribed
standard part. Their difference is multiplied, in the Taylor difference
identity, by a matrix whose standard part is the invertible matrix
$A(\st x)$. That matrix is invertible over $\K$, since its determinant
has nonzero standard part. Therefore the difference is zero.
\end{proof}
The theorem explains why analytic systems with simple residual Jacobian
are stable under arbitrarily deep set-supported surreal perturbations.
It does not require an ordinary convergent iteration in $\K$.

\section{Worked examples and constructive calculations}
\subsection{An exact infinitesimal Lambert-type solution}
Let $\eta\in\m$. The coherent equation
\[
 z-\eta\widetilde{e^z}=0
\]
has exactly one root in the finite tube: its leading coefficient is $z$,
so all finite roots have standard part zero and the zero-cluster count
there is one. Lagrange inversion gives
\begin{equation}\label{eq:tree}
 z=\sum_{n\geq1}\frac{n^{n-1}}{n!}\eta^n
 =\eta+\eta^2+\frac32\eta^3+\frac83\eta^4
     +\frac{125}{24}\eta^5+\frac{54}{5}\eta^6+\cdots.
\end{equation}
The sum is valid for \emph{every} surcomplex infinitesimal $\eta$, regardless
of the complexity of its positive support. It solves the equation
exactly by formal substitution and summability. This makes no claim about
possible infinite roots of a globally chosen exponential equation.

\subsection{First preparation coefficients for a double root}
For $A(w)=w^2-t e^w$, where $t=\omega^{-1}$, take $m=2$. Put
\[
 B(w)=\frac{e^w-1-w}{w^2},
\]
with its removable value at zero. The first homogeneous terms in
\cref{thm:preparation} are
\[
 p^{[1]}=-t(1+w),\qquad Q^{[1]}=-tB(w),
\]
and
\[
 p^{[2]}=-t^2\left(\frac12+\frac23w\right).
\]
Indeed, $p^{[1]}Q^{[1]}=t^2(1+w)B(w)$, whose constant and linear Taylor
coefficients are $1/2$ and $2/3$. Therefore the prepared polynomial begins
\[
 P(w)=w^2-t(1+w)-t^2\left(\frac12+\frac23w\right)+\cdots.
\]
Its two roots are precisely the two roots of $w^2-te^w$ in $\m$,
counted with multiplicity. The preparation recursion computes them without
assuming an a priori square-root splitting of the perturbation.

\subsection{Two genuinely separated surreal scales}
Let
\[
 \lambda=\omega^{-1/2},\qquad s=\omega^{-\omega},\qquad
 F(z)=z^2-\lambda^2-s\widetilde{e^z}.
\]
The perturbation $s$ is smaller than every positive ordinary power of
$\lambda$. The leading coefficient $z^2$ gives exactly two roots in
$\m$. Set $z=\lambda w$ and divide by $\lambda^2$. The normalized leading
coefficient is $w^2-1$, so one root lies near each of $w=1,-1$.

Writing $r=\lambda$ or $r=-\lambda$, substitution gives
\begin{equation}\label{eq:two-scale}
 z_r=r+s\frac{e^r}{2r}
       +s^2\frac{e^{2r}(2r-1)}{8r^3}
       \pmod {s^3}.
\end{equation}
Here $e^r$ is the infinitesimal Taylor exponential. The congruence is in
$\C((\lambda))[[s]]$, embedded in $\K$ by the ordered exponents
$n\omega+k/2$. It is an $s$-adic statement, not an assertion that the
remainder is bounded by a fixed real multiple of $s^3$.

To verify the displayed coefficients, put $z=r+sa+s^2b$ and use
$r^2=\lambda^2$. The coefficients of $s$ and $s^2$ in $F(z)$ are
\[
 2ra-e^r,\qquad a^2+2rb-e^ra.
\]
They vanish exactly for $a=e^r/(2r)$ and
$b=e^{2r}(2r-1)/(8r^3)$. Higher coefficients are determined similarly.

\subsection{Arbitrary well-ordered families of perturbation scales}
Let $E\subset\No^{>0}$ be any well-ordered set, choose arbitrary ordinary
complex constants $c_\gamma$, and let $m\geq1$. Then
\[
 F(z)=z^m+\sum_{\gamma\in E}t^\gamma e^{c_\gamma z}
\]
is a coherent section on the finite plane and has exactly $m$ roots in
$\m$, counted with multiplicity. No boundedness condition on the family
$c_\gamma$ is needed. The positive support is essential: a decreasing
sequence of exponents tending to zero would not be a permitted Hahn
support. For example $\{1/n:n\geq1\}$ is excluded, whereas the increasing
support $\{1-1/n:n\geq2\}$ is permitted.

\subsection{What can be computed finitely}
Preparation and reversion are constructive when the input supplies
coefficient functions and an effective representation of the relevant
supports. A finite computation can be organized by positive-support word
degree. To compute degree $N$ in preparation, use
\eqref{eq:preprecursion} for $1\leq n\leq N$; only finite products and
Taylor division operators occur at each degree. With finitely many formal
perturbation generators, the coefficient at degree $n$ is a finite sum of
analytic expressions associated with binary trees.

For a completely arbitrary surreal support, the existence theorem is not
a promise that all coefficients can be enumerated by a finite program.
The accompanying \texttt{verify.py} instead tests the displayed finite
algebraic consequences with exact rational arithmetic. It checks the
Lambert-type solution through degree forty, the preparation identity in
a finite bivariate truncation, and the two-scale root residual. These are
reproducible checks of formulas, not substitutes for the support proofs.

\Needspace{14\baselineskip}
\section{What has been established, and what remains to be chosen}
\subsection{A comparison of the analytic categories}
\begin{longtable}{>{\raggedright\arraybackslash}p{.22\textwidth}>{\raggedright\arraybackslash}p{.32\textwidth}>{\raggedright\arraybackslash}p{.37\textwidth}}
\toprule
\textbf{Category} & \textbf{Defining structure} & \textbf{Conclusions and restrictions}\\
\midrule
\endhead
Differential holomorphy & Full-surreal $\eta$--$\delta$ complex derivative &
Product and chain rules; locally constant counterexamples prevent general
identity and Liouville theorems.\\[4pt]
Monadic formal functions & Evaluation of ordinary formal series at
infinitesimals & Taylor calculus includes complex-divergent series;
no coherence between different standard parts is imposed.\\[4pt]
Canonical lifts & Taylor extension of one ordinary holomorphic function &
Classical analytic continuation and sharp lifted Schwarz lemma;
only the specified ordinary tube is covered.\\[4pt]
Hahn-coherent sections & A common ordinary domain and a set-sized
well-ordered scale support & Cauchy and residue calculus, preparation,
exact zero counts, open mapping and local inverse theorems.\\[4pt]
All-scale coherence & A Hahn-coherent representative after every positive
surreal dilation & Exactly the finite-degree polynomials; a global
Liouville theorem follows.\\[4pt]
Canonical global exponential & Real surreal exponential plus a normal-form
phase convention & A surjective exponential with kernel $2\pi i\Oz$;
locally coherent but not all-scale coherent.\\
\bottomrule
\end{longtable}

\subsection{Natural extensions of the framework}
The results isolate several precise directions rather than one vague
request for ``analysis on all surreals.'' One can enlarge the coefficient
category from holomorphic functions to specified asymptotic analytic
families and ask which support conditions survive changes by infinite
normalized scales. One can compare such a category with definable complex
analysis, where definable connectedness supplies a different global
coherence mechanism. One can also study divisors and ramification in
families of Hahn sections using the multivariable implicit theorem and
local preparation.

A broader transcendental global theory must relax the all-scale Hahn
condition: \cref{thm:rigidity} proves that merely insisting on more charts
of the same type cannot achieve that goal. For exponentials, one must
also specify which purely infinite phases are identified. The explicit
characters in \cref{thm:twists} show that the local equation $E'=E$ cannot
make that choice. These are structural design constraints established by
theorems, not gaps to be hidden by unqualified uses of ``holomorphic.''

\subsection{A practical proof checklist}
A proposed surcomplex analytic argument should identify its coefficient
field, its set support, its summability convention, the topology in any
limit, and the coherence connecting its local formulas. In a contour
argument it should additionally identify whether integration is ordinary,
coefficientwise, definable, or a separately constructed surreal integral.
When counting zeros, it should state whether multiplicities refer to
ordinary leading coefficients, actual surcomplex roots, or formal
polynomial factors. The theorems above explicitly connect all three
counts in the prepared setting.

\section*{Conclusion}
\addcontentsline{toc}{section}{Conclusion}
The field $\No[i]$ supports more than a formal analogy with $\C$.
Hahn-coherent sections admit an actual differential calculus, exact
infinitesimal root counts, local mapping theorems, and a coefficientwise
Cauchy--residue theory. At the same time, the proper-class abundance of
scales makes ordinary sequential arguments unavailable and turns
all-scale coherence into a polynomial rigidity condition. A successful
surcomplex analysis therefore consists of explicit compatible structures,
not a single unqualified substitution of fields. The framework developed
here provides such structures and proves both their strongest elementary
consequences and the boundaries separating them.

\appendix
\section{A formal Lagrange coefficient calculation}\label{app:lagrange}
Let $u=q\phi(u)$, where $q$ is an ordinary formal parameter and
$\phi(X)$ is a formal series over a characteristic-zero coefficient ring.
The recursive solution $u\in qA[[q]]$ is unique. For $n\geq1$,
\begin{equation}\label{eq:lagrange}
 [q^n]u=\frac1n[X^{n-1}]\phi(X)^n.
\end{equation}
Here is the formal residue proof when $\phi(0)$ is invertible. Change
variables $q=X/\phi(X)$ and use \cref{thm:reschange}:
\begin{align*}
 [q^n]u
 &=\Res_{q=0}u\,q^{-n-1}\dd q\\
 &=\Res_{X=0}\left(\phi(X)^nX^{-n}
       -\phi(X)^{n-1}\phi'(X)X^{1-n}\right)\dd X\\
 &=\left(1-\frac{n-1}{n}\right)[X^{n-1}]\phi(X)^n.
\end{align*}
For each fixed $n$, both sides are universal polynomials with rational
coefficients in the first $n$ coefficients of $\phi$. The identity in
the localization where the constant coefficient is invertible therefore
holds in the universal polynomial ring and then under every
specialization. This removes the invertibility assumption.

For reversion take $\phi(X)=-H(y+X)$ to obtain
\eqref{eq:reversion}. For \eqref{eq:tree}, take $\phi(X)=e^X$:
\[
 \frac1n[X^{n-1}]e^{nX}=\frac{n^{n-1}}{n!}.
\]
Only after these ordinary formal coefficient identities are established
do we substitute positive-support Hahn parameters.

\section{Foundational inputs and verification boundaries}
The arguments use the following established inputs: the surreal cut
property and normal-form theorem; real-closedness of $\No$; Neumann's
support lemma; the group law and infinitesimal Taylor behavior of the real
surreal exponential; and ordinary complex analytic theorems on the stated
ordinary domains. Preparation, the exact zero-cluster count, the
full-surreal Taylor estimate, the local mapping reductions, all-scale
rigidity, the explicit exponential twists, and the stated value-distribution
calculation are proved in the article from those inputs.

The bibliography distinguishes source literature from the present proofs.
In particular, the existence and kernel of the canonical surcomplex
exponential are attributed to Ehrlich and Kaplan, and the general idea of
surreal analytic functions is attributed to the much earlier work of
Alling. The finite checks supplied with the article have been executed,
but neither the full article nor the foundational inputs have been
machine-checked in a proof assistant. A literature search cannot certify
priority for every formulation, and none is asserted here.

\begin{thebibliography}{99}
\raggedright
\bibitem{Conway}
J. H. Conway, \emph{On Numbers and Games}, second edition,
A K Peters, 2001; first edition, Academic Press, 1976.

\bibitem{Gonshor}
H. Gonshor, \emph{An Introduction to the Theory of Surreal Numbers},
London Mathematical Society Lecture Note Series 110,
Cambridge University Press, 1986.
\href{https://doi.org/10.1017/CBO9780511629143}{doi:10.1017/CBO9780511629143}.

\bibitem{Alling}
N. L. Alling, \emph{Foundations of Analysis over Surreal Number Fields},
North-Holland Mathematics Studies 141, North-Holland, 1987.
\href{https://shop.elsevier.com/books/foundations-of-analysis-over-surreal-number-fields/alling/978-0-444-70226-5}{Publisher's description and bibliographic record}.

\bibitem{AllingShort}
N. L. Alling, Fundamentals of analysis over surreal number fields,
\emph{Rocky Mountain Journal of Mathematics} \textbf{19} (1989), 565--573.
\href{https://doi.org/10.1216/RMJ-1989-19-3-565}{doi:10.1216/RMJ-1989-19-3-565}.

\bibitem{Neumann}
B. H. Neumann, On ordered division rings,
\emph{Transactions of the American Mathematical Society}
\textbf{66} (1949), 202--252.
\href{https://doi.org/10.1090/S0002-9947-1949-0032593-5}{doi:10.1090/S0002-9947-1949-0032593-5}.

\bibitem{vdDE}
L. van den Dries and P. Ehrlich, Fields of surreal numbers and
exponentiation, \emph{Fundamenta Mathematicae} \textbf{167} (2001), 173--188.
\href{https://doi.org/10.4064/fm167-2-3}{doi:10.4064/fm167-2-3}.

\bibitem{EK}
P. Ehrlich and E. Kaplan, Surreal ordered exponential fields,
\emph{The Journal of Symbolic Logic} \textbf{86} (2021), 1066--1115.
\href{https://doi.org/10.1017/jsl.2021.59}{doi:10.1017/jsl.2021.59}.
Author preprint: \href{https://arxiv.org/abs/2002.07739}{arXiv:2002.07739v3},
especially Section 11. See also the reference-correction erratum,
\href{https://doi.org/10.1017/jsl.2022.5}{doi:10.1017/jsl.2022.5}.

\bibitem{RSS}
S. Rubinstein-Salzedo and A. Swaminathan, Analysis on surreal numbers,
\emph{Journal of Logic \& Analysis} \textbf{6}:5 (2014), 1--39.
\href{https://arxiv.org/abs/1307.7392}{arXiv:1307.7392}.

\bibitem{PS}
Y. Peterzil and S. Starchenko, Expansions of algebraically closed fields
in o-minimal structures, \emph{Selecta Mathematica, New Series}
\textbf{7} (2001), 409--445.
\href{https://doi.org/10.1007/PL00001405}{doi:10.1007/PL00001405}.
See also the authors'
\href{https://math.haifa.ac.il/kobi/Correction.pdf}{correction}.

\bibitem{PSsurvey}
Y. Peterzil and S. Starchenko, Complex analytic geometry in a nonstandard
setting, in \emph{Model Theory with Applications to Algebra and Analysis},
Cambridge University Press, 2008, 117--166.
\href{https://math.haifa.ac.il/kobi/Newton.pdf}{Author manuscript}.

\bibitem{Orloff}
J. Orloff, \emph{Complex Variables with Applications}, MIT course 18.04,
Spring 2018, lecture notes. Background for ordinary Cauchy theory,
residues, the argument principle, and related classical results.
\href{https://ocw.mit.edu/courses/18-04-complex-variables-with-applications-spring-2018/pages/lecture-notes/}{MIT OpenCourseWare lecture notes}.

\bibitem{MantovaMatusinski}
V. Mantova and M. Matusinski, Surreal numbers with derivation,
Hardy fields and transseries: a survey,
\href{https://arxiv.org/abs/1608.03413}{arXiv:1608.03413}, 2016.
Background on normal forms, summability, and the distinction between
function differentiation and derivations of surreal fields.
\end{thebibliography}
\end{document}
